\problem{1}
\textbf{[Отбор 2024, 7.1]}
Барон Мюнхгаузен утверждает, что смог нарисовать на плоскости семь отрезков так, что каждый пересекается ровно с четырьмя другими. Докажите, что ему можно верить.

\problem{2}
\textbf{[Отбор 2024, 7.3]} 
В равнобедренном треугольнике $A B C(A B=B C)$ на сторонах $A C$ и $A B$ выбраны точки $P$ и $Q$ соответственно. Оказалось, что $A P=A B$ и $P Q=P B$. Докажите, что $A Q=C P$.

\problem{3}
\textbf{[Отбор 2023, 7.3]}
Дан треугольник $A B C(A B>B C)$. Через точку $A$ провели луч параллельный $B C$, а через точку $C$ - луч параллельный $A B$. Эти лучи пересеклись в точке $D$. На луче $B D$ выбраны различные точки $X$ и $Y$ так, что $B C=C X$ и $B A=A Y$. Докажите, что $X D=D Y$.

\problem{4}
\textbf{[Отбор 2022, 6.4]}
Отрезок, длина которого 88 cm . разделён на четыре неравных отрезка. Расстояние между серединами самого левого и самого правого отрезков равно 46 cm . Найдите расстояние между серединами средних отрезков.

\problem{5}
\textbf{[Отбор 2022, 7.5]}
В треугольнике $A B C$ угол $A$ в 3 раза больше угла $B$, который в 2 раза больше угла $C$. Сторона $B C$ на 5 длиннее стороны $A C$. Найдите длину биссектрисы треугольника, проведенной из вершины $A$.

\problem{6}
\textbf{[Отбор 2021, 7.2]}
В остроугольном треугольнике $A B C$ точки $D$ и $E$ - основания высот, проведенных из вершин $B$ и $C$ соответственно. $\angle C B D=2 \angle A B D$ и $\angle A C E=3 \angle B C E$. Найдите углы треугольника $A B C$.

\problem{7}
\textbf{[Отбор 2021, 7.4]}
В выпуклом четырехугольнике $A B C D$ равны углы $A, B, C$. Кроме того равны стороны $A D$ и $C D$. На стороне $A B$ выбрана такая точка $E$, что $B E=A D$. Докажите, что $C E$ - биссектриса угла $B C D$.